\section{Introduction}
\label{firstchapter}

The environmental impact of structural systems has become a central consideration in contemporary building design, particularly in light of increasing regulatory and societal pressure to reduce greenhouse gas emissions. The construction sector is a major contributor to global carbon emissions, and demand for built space continues to grow: the required usable floor area worldwide is projected to increase substantially by 2060, driven by rising living standards, urbanisation, population growth, and the replacement of ageing infrastructure. \todo[inline]{Add citation for the floor-area projection / sustainable development source referenced in the original draft, and confirm the 2060 figure.}

Within multi-storey buildings, floor slabs typically account for a substantial proportion of total material consumption and thus represent a significant contributor to embodied carbon. Because slabs are repeated on every storey and span large surface areas, even modest reductions in material per square metre accumulate into significant savings across a building. The choice of floor system therefore has a significant influence on the embodied carbon of the overall structure, making it a high-value target for environmental optimisation.

While reinforced concrete (RC) flat slabs remain the predominant solution in many regions, composite slim floor systems offer an alternative approach that may reduce material usage and overall Global Warming Potential (GWP). In the United Kingdom in particular, composite floor systems dominate the multi-storey steel-framed building market \citep{RefWorks:2011composite}, supported by mature design guidance and an established supply chain of profiled steel decking products.

Improving structural efficiency is an important strategy for reducing the environmental impact of buildings. Structural efficiency can be understood as the ability of a structural system to carry the required loads while minimising material consumption. One approach to improving structural efficiency is the use of composite construction, in which different materials are combined to exploit their respective mechanical properties. Steel--concrete composite systems take advantage of the compressive strength of concrete and the tensile capacity of steel, allowing structural elements to achieve higher load-bearing efficiency than conventional reinforced concrete systems. Slim floor systems extend this principle further by integrating the steel beam within the depth of the slab, reducing overall structural depth and, by extension, the floor-to-floor height of the building.

Despite these potential advantages, the relationship between structural efficiency and environmental performance is not always straightforward. Structural steel generally has a significantly higher embodied carbon per unit mass than concrete. As a result, a system that reduces concrete volume through the introduction of structural steel may not necessarily achieve a lower overall GWP. The environmental performance of such systems therefore depends on the balance between reduced concrete usage and increased steel consumption. Quantitative assessment is required to determine whether alternative structural systems provide genuine environmental benefits.

This tension is particularly relevant for composite slabs, where the profiled steel deck serves a dual role: it acts as permanent formwork and as tensile reinforcement in the completed slab. Whether the structural and constructional advantages of this arrangement translate into a net reduction in embodied carbon, relative to a conventional RC flat slab designed to the same requirements, is the central question this thesis addresses.

This study builds on a pre-existing parametric optimisation framework developed at the Chair of Concrete Structures and Bridge Design at ETH Zurich \todo[inline]{Add citation for the prior ETH Zurich framework --- ask A+T for the correct reference}. The original framework provides SIA-compliant design checks and GWP quantification for reinforced concrete and timber slab systems, identifying minimum-carbon cross-sections across a range of spans. However, it does not currently include composite steel--concrete slab systems, leaving a gap in the comparison of structural typologies on an embodied-carbon basis.

The principal contribution of this thesis is the extension of this framework to incorporate composite slab cross-sections. A new \texttt{CompositeSlab} class and associated optimisation routines were developed, supported by a database of commercially available profiled steel decking products with their mechanical properties and EPD-derived GWP values. This enables, for the first time within the framework, a systematic and consistent comparison between composite slabs and the existing reinforced concrete baselines, using identical loading assumptions, design standards, and environmental accounting. The composite slab model implements the full staged design of composite construction --- construction stage, composite stage, fire, and serviceability checks --- in accordance with the Eurocodes, so that all candidate cross-sections satisfy current design requirements before their GWP is compared.

The overall aim of this thesis is to evaluate the potential of composite steel--concrete floor slabs for sustainable structural design, by quantifying their embodied carbon relative to conventional reinforced concrete slabs across a range of practical spans and occupancy scenarios.

\todo[inline]{Write this subsection once the results are finalised. Suggested content: (1) the governing limit state for composite slabs across the span range (SLS deflection dominant, fire D1 at short spans); (2) the headline comparison against the RC baselines and the crossover span (if any); (3) any notable scenario-dependent differences (office vs retail vs industrial). Keep to one short paragraph --- this previews the conclusions without repeating Chapter~\ref{fourthchapter}.}


The remainder of this thesis is structured as follows. Chapter~\ref{secondchapter} reviews the relevant literature, covering sustainability in structural design, the development and typologies of composite and slim floor systems, the Eurocode design methods for composite slabs, and existing optimisation studies, concluding with the research gap addressed by this work. Chapter~\ref{thirdchapter} sets out the methodology, including the pre-existing framework, the material and product database, the case studies and loading scenarios, the structural design procedure, the environmental impact calculation, the optimisation algorithm, and the validation against manufacturer data. Chapter~\ref{fourthchapter} presents the optimisation results, including the validation outcomes, the minimum-GWP cross-sections, the governing limit states, and the GWP breakdown by material. Chapter~\ref{fifthchapter} discusses the implications of these results, and Chapter~\ref{sixthchapter} draws conclusions and outlines recommendations for further work.
